% 库拉托夫斯基闭包公理（Kuratowski closure axioms）（综述）
% license CCBYSA3
% type Wiki

本文根据 CC-BY-SA 协议转载翻译自维基百科\href{https://en.wikipedia.org/wiki/Kuratowski_closure_axioms}{相关文章}。

在拓扑学及其相关数学分支中，库拉托夫斯基闭包公理（Kuratowski closure axioms）是一组可用于在一个集合上定义拓扑结构的公理。它们与更常用的开集定义是等价的。该体系最早由卡齐米日·库拉托夫斯基（Kazimierz Kuratowski）形式化提出，\(^\text{[1]}\)后来又被数学家如瓦茨瓦夫·谢尔平斯基（Wacław Sierpiński）和安东尼奥·蒙泰罗（António Monteiro）等人进一步研究。\(^\text{[2]}\)

另一组类似的公理也可以仅通过其对偶概念——内点算子（interior operator）来定义拓扑结构。\(^\text{[3]}\)
\subsection{定义}
\subsubsection{库拉托夫斯基闭包算子及其弱化形式}
设$X$为任意集合，$\wp(X)$为其幂集。一个库拉托夫斯基闭包算子（Kuratowski closure operator）是一个一元运算$\mathbf{c} : \wp(X) \to \wp(X)$满足以下性质：

\textbf{[K1]}保持空集：$\mathbf{c}(\varnothing) = \varnothing
$

\textbf{[K2]}外延性（Extensive）：对所有$A \subseteq X $，$A \subseteq \mathbf{c}(A)$

\textbf{[K3]}幂等性（Idempotent）：对所有$A \subseteq X$，$\mathbf{c}(A) = \mathbf{c}(\mathbf{c}(A))$

\textbf{[K4]}对二元并集保持（或分配律）：对所有$A, B \subseteq X$，
$\mathbf{c}(A \cup B) = \mathbf{c}(A) \cup \mathbf{c}(B)$

由于闭包算子$\mathbf{c}$保持二元并集，上述条件的一个推论是\(^\text{[4]}\)：

\textbf{[K4′]}单调性（Monotone）：$A \subseteq B \Rightarrow \mathbf{c}(A) \subseteq \mathbf{c}(B)$

事实上，如果我们将\textbf{[K4]}中的等号改为包含号，就得到一个更弱的公理\textbf{[K4″]}（次可加性，subadditivity）：

\textbf{[K4″]}次可加性：对所有$A, B \subseteq X $，$\mathbf{c}(A \cup B) \subseteq \mathbf{c}(A) \cup \mathbf{c}(B)$

那么可以容易看出，公理\textbf{[K4′]}与\textbf{[K4″]}一起等价于 \textbf{[K4]}（见下文“证明2”的倒数第二段）。

库拉托夫斯基（1966）提出了一个第五个（可选）公理，要求单点集在闭包下保持不变：对所有$ x \in X $，$\mathbf{c}({x}) = {x}$他将满足全部五条公理的拓扑空间称为T₁ 空间（T1-spaces），以区别于仅满足前四条公理的一般空间。实际上，这些空间与通常意义下的拓扑T₁ 空间完全对应（见下文）。\(^\text{[5]}\)

如果省略条件\textbf{[K3]}，则这些公理定义了一个 Čech 闭包算子。\(^\text{[6]}\)若改为省略\textbf{[K1]}，则满足\textbf{[K2]}、\textbf{[K3]}与 \textbf{[K4′]}的算子称为 Moore 闭包算子。\(^\text{[7]}\)因此，一个二元组$(X, \mathbf{c})$称为库拉托夫斯基空间（Kuratowski closure space）、Čech 闭包空间（Čech closure space）或Moore 闭包空间（Moore closure space），取决于闭包算子$\mathbf{c}$满足的公理体系。
\subsubsection{替代公理化}
四条库拉托夫斯基闭包公理可以被一个条件所取代，这一条件由 Pervin 提出：[8]

$\text{[P]}\quad \forall, A,B \subseteq X, \quad A \cup \mathbf{c}(A) \cup \mathbf{c}(\mathbf{c}(B)) = \mathbf{c}(A \cup B) \setminus \mathbf{c}(\varnothing)$
即：
对所有$A, B \subseteq X$，有$A \cup c(A) \cup c(c(B)) = c(A \cup B) \setminus c(\varnothing)$

由该条件可以推出公理\textbf{[K1]–[K4]}：
\begin{enumerate}
\item 取$A = B = \varnothing$。则$\varnothing \cup c(\varnothing) \cup c(c(\varnothing)) = c(\varnothing) \setminus c(\varnothing) = \varnothing$即$c(\varnothing) \cup c(c(\varnothing)) = \varnothing$这立即推出\textbf{[K1]}。
\item 取任意$A \subseteq X$，令$B = \varnothing$。应用\textbf{[K1]}，得$A \cup c(A) = c(A)$推出\textbf{[K2]}。
\item 取$A = \varnothing$，任意$B \subseteq X$。应用\textbf{[K1]}，得$c(c(B)) = c(B)$即\textbf{[K3]}。
\item 取任意 (A, B \subseteq X)，结合 [K1]–[K3]，可推得 **\textbf{[K4]}。
\end{enumerate}
另一种替代方案由 Monteiro（1945）提出，他提出了一个更弱的公理，仅能推出 \textbf{[K2]–[K4]}：[9]

$\text{[M]}\quad \forall, A,B \subseteq X, \quad A \cup c(A) \cup c(c(B)) \subseteq c(A \cup B)$

即：$A \cup c(A) \cup c(c(B)) \subseteq c(A \cup B)$

条件\textbf{[K1]}与\textbf{[M]}相互独立。确实，如果$X \neq \varnothing$，定义算子$\mathbf{c}^\star : \wp(X) \to \wp(X)$为常值映射$A \mapsto \mathbf{c}^\star(A) := X$则该算子满足\textbf{[M]}，但不保持空集，因为$\mathbf{c}^\star(\varnothing) = X$由定义可见，任何满足 \textbf{[M]} 的算子都是一个 Moore 闭包算子。

M. O. Botelho 与 M. H. Teixeira** 提出了一个更对称的替代形\textbf{[BT]}，它同样蕴含公理\textbf{[K2]–[K4]}：[2]

$\text{[BT]}\quad \forall, A,B \subseteq X, \quad A \cup B \cup c(c(A)) \cup c(c(B)) = c(A \cup B)$

即：$A \cup B \cup c(c(A)) \cup c(c(B)) = c(A \cup B)$

